\section{Forecasting — Series de Tiempo}
\label{sec:forecasting}

\subsection*{Planteamiento}

Se busca pronosticar la energía semanal generada por la planta, información útil para planeación comercial y operativa (p. ej. compromisos de entrega de energía, programación de mantenimiento en semanas de baja generación esperada).

La serie de 5 minutos se agregó a frecuencia \textbf{semanal} (suma de \texttt{total\_ac\_power}), reduciendo el ruido de alta frecuencia y haciendo tratable la estacionalidad anual (\texttt{seasonal\_periods=52}). Esto produjo una serie de \textbf{313 semanas} (2017--2023).

\begin{figure}[H]
\centering
\includegraphics[width=\textwidth]{fig_serie_semanal.png}
\caption{Energía AC total generada por semana (2017--2023).}
\end{figure}

\subsection*{Descomposición estacional}

\begin{figure}[H]
\centering
\includegraphics[width=0.85\textwidth]{fig_descomposicion_estacional.png}
\caption{Descomposición aditiva de la serie semanal: tendencia, estacionalidad anual y residuo.}
\end{figure}

La descomposición confirma una estacionalidad anual marcada y regular (picos y valles que se repiten en la misma época cada año, asociados a la posición del sol) superpuesta a una tendencia de más largo plazo; el componente residual conserva variabilidad no explicada por tendencia ni estacionalidad, atribuible a condiciones climáticas semana a semana (nubosidad, lluvia) no registradas en el dataset.

\subsection*{Modelos comparados}

\begin{itemize}
    \item \textbf{Baseline ingenuo estacional}: el valor de la misma semana del año anterior.
    \item \textbf{Holt-Winters} (suavizado exponencial triple) \cite{holt_winters}: captura tendencia y estacionalidad anual de forma explícita.
\end{itemize}

Se reservó el último año (52 semanas) como conjunto de prueba, entrenando ambos modelos únicamente con los datos anteriores.

\begin{table}[H]
\centering
\begin{tabular}{lrrr}
\toprule
\textbf{Modelo} & \textbf{RMSE} & \textbf{MAE} & \textbf{MAPE (\%)} \\
\midrule
Naive estacional & 110\,050.42 & 81\,013.18 & 29.64 \\
Holt-Winters      & 66\,010.72  & 52\,908.83  & 21.79 \\
\bottomrule
\end{tabular}
\caption{Comparación de modelos de forecasting sobre las últimas 52 semanas (conjunto de prueba)}
\end{table}

\begin{figure}[H]
\centering
\includegraphics[width=\textwidth]{fig_forecast.png}
\caption{Pronóstico semanal — real vs. baseline ingenuo vs. Holt-Winters.}
\end{figure}

\textbf{Modelo seleccionado:} Holt-Winters, con una reducción de $\sim$40\% en RMSE y MAE, y de casi 8 puntos porcentuales de MAPE (21.79\% vs. 29.64\%), respecto al baseline ingenuo estacional. Esto confirma que modelar explícitamente la tendencia además de la estacionalidad aporta valor frente a simplemente repetir el valor del año anterior.

\subsection*{Significado práctico y limitaciones}

Un pronóstico semanal con error del orden de 21.8\% (MAPE) permite anticipar, con semanas de anticipación, si la producción estará por debajo de lo esperado y planear en consecuencia. La principal limitación es que el modelo es \textbf{univariado}: no incorpora variables meteorológicas externas (irradiancia, nubosidad, temperatura ambiente) que explican gran parte de la variabilidad semana a semana, ni eventos de mantenimiento programado; por ello su desempeño se degrada en semanas con condiciones climáticas atípicas respecto al patrón histórico.
